% Page 019 — page imprimée 15
% Une épidémie à Meyrueis en 1768
% Reconstruction éditoriale en LaTeX.

{\fontsize{10.25}{12.1}\selectfont

\begin{center}
  {\bfseries\large Une épidémie à Meyrueis en 1768}
\end{center}

\vspace{0.55em}

\begin{center}
\fbox{%
\begin{minipage}{0.91\textwidth}
\itshape
Ce compte rendu d’inspection médicale à Meyrueis et dans les environs a été rédigé en 1768, vingt ans avant la révolution. On peut le consulter au Service du Patrimoine de la Médiathèque de Montpellier sous la cote 19.180.

\vspace{0.35em}

Il m’a vivement intéressé pour trois raisons :

\vspace{0.15em}

1/ \quad La description assez précise de cette épidémie dénommée « peste blanche » conséquence du manque d’hygiène, de la pauvreté et d’une nourriture insuffisante et souvent avariée.

\vspace{0.15em}

2/ \quad L’évocation de la vie quotidienne des paysans de Meyrueis et des vallées des environs à la suite de mauvaises récoltes dues à des sécheresses persistante et à des hivers rigoureux. Les maisons ne sont que des masures sans aération où « l’on meurt de froid ou de chaud », la promiscuité est générale, excellent terreau pour la tuberculose et toutes les maladies infectieuses.

\vspace{0.15em}

3/ \quad Le style et les tournures surannées de ces lignes.

\vspace{0.15em}

Ce mémoire complète donc ce que l’on peut savoir de la vie de nos aïeux à Meyrueis mais aussi à Carnac près de Rousses où la famille Couderc résidait alors.
\end{minipage}}
\end{center}

\vspace{0.45em}

\begin{center}
{\fontsize{18}{20}\selectfont\bfseries MEMOIRE\par}
\vspace{0.55em}
{\bfseries
sur la maladie épidémique\\
qui a régné à Meyrueis et ses environs
}
\vspace{0.15em}

Présenté à Monsieur de Saint Priest\\
Le 3 Septembre 1768

\vspace{0.35em}

{\bfseries Par Tandon, Docteur en Médecine\par}
Imprimé par Auguste François Rochard, Place du petit Scel\\
Montpellier 1769
\end{center}

\vspace{0.55em}

C’est une maladie qui a régné assez longtemps dans une partie très considérable des Hautes Cévennes, manifestée par des fièvres intermittentes d’un très mauvais caractère. Elles ont d’abord sévi à St Jean de Védas, la Verune, Fabrègues, Gigean, Loupian, Cournon venant des exhalaisons des étangs par fortes chaleurs et de mauvaises récoltes. Le plus grand nombre des habitants est obligé de se nourrir d’\textit{azéroles*}, de mûres de haies, de figues et de raisins qui n’étaient pas mûrs et de petites carpes, muges, anguilles et loups.

Puis l’épidémie s’est répandue sur Marsillargue, Clermont, Lodève, Mèze et enfin les Hautes Cévennes.

\vspace{0.3em}

{\itshape\textbf{Azérole :} Fruit de l’azérolier, jaune ou rouge, ressemblant à une petite pomme}

}

